\textbf{\textit{Resumen}} — La seguridad en las comunicaciones de drones de reparto constituye un desafío crítico en la logística moderna. Los enlaces telemáticos entre drones y estaciones base son vulnerables a ataques de suplantación de identidad, interceptación, inyección de comandos y repetición. Este trabajo presenta el diseño, implementación y validación de un esquema de cifrado híbrido de cuatro capas que integra una \textbf{PUF Neuromórfica} con simulación de \textbf{memristor crossbar} y plasticidad sináptica dependiente de disparos (STDP) como generador de material clave raíz. El sistema emplea una cascada de cifrados modernos: \textbf{AES-256-CBC} con relleno PKCS7 como primera capa de confusión, \textbf{AES-256-CTR} como cifrado de flujo intermedio, \textbf{Permutación Spike} basada en el patrón de disparos neuronales como capa de difusión no lineal, y \textbf{ChaCha20} como cifrado de flujo final. El protocolo de autenticación se fundamenta en pares Reto-Respuesta (CRP) generados por la PUF neuromórfica, los cuales son verificados contra una base de datos PostgreSQL que almacena exclusivamente los CRP ya utilizados (helper data), nunca la clave maestra. Se implementó un prototipo funcional en Python 3.13 con simulación de vuelo completa, generación de telemetría cifrada, y detección de drones falsos. Los resultados experimentales sobre 1,000 iteraciones demuestran: (1) 100\% de éxito en cifrado y descifrado; (2) detección de suplantación con probabilidad superior a $1 - 5.4 \times 10^{-20}$; (3) latencia media de cifrado de 2.1 ms por paquete; (4) evolución dinámica de las conductancias del crossbar memristivo mediante aprendizaje STDP, adaptando la función PUF tras cada autenticación exitosa. El trabajo concluye que la arquitectura neuromórfica es viable para sistemas embebidos y sienta las bases para implementaciones con PUF memristivas reales en silicio.
\vspace{0.5em}

\textbf{\textit{Index Terms}} — PUF, Physical Unclonable Function, Neuromorphic Computing, Memristor Crossbar, STDP, AES-256, Cifrado Híbrido, Dron de Reparto, Seguridad Telemática, ChaCha20, Autenticación Hardware, CRP.
